% Generated by tables/make_methods-compared.py from data/derived/benchmark/methods-compared.json; do not edit.
\subsubsection{Reverse liquidity stress}
\label{sec:results-methods-compared-q10}

This question is answered on case C16. One bank, one balance sheet and one labelled run: every variant document scores the modelled outflow against the same observed share of deposits lost in the two days of March 2023. The ratio and the light are read against that observed share, which is a label and never a forecast. The run-off multiplier is a threshold on a base, not a forecast: it has no score of its own, so it is shown beside the others and not ranked. The run ``The labelled run'' reads as follows: Balance sheet at 2022-12-31 with synthetic deposit behaviour; event of 9 and 10 March 2023, horizon 2 days; its random seed is 42. Table~\ref{tab:bm-q10-methods} lists the five methods the document compares, one long-established, two recent published, one proposed in this work and one reference, with the section of this report that describes each. Tables~\ref{tab:bm-q10-1}, \ref{tab:bm-q10-2} give the values, one table for each group of metrics that share a run and a layout.

{\footnotesize\setlength{\tabcolsep}{3pt}
% bm-table-methods Q10
\begin{longtable}{>{\raggedright\arraybackslash}p{0.38\textwidth}>{\raggedright\arraybackslash}p{0.17\textwidth}>{\raggedright\arraybackslash}p{0.27\textwidth}>{\raggedright\arraybackslash}p{0.10\textwidth}}
\caption{The methods compared on the question Reverse liquidity stress: the standing of each method, the source and the year the document records for it, and the section of this report that describes it; a dash marks a method with no source recorded or with no section of its own.}\label{tab:bm-q10-methods}\\
\toprule Method & Standing & Source & Section \\ \midrule \endfirsthead
\toprule Method & Standing & Source & Section \\ \midrule \endhead
Liquidity coverage ratio formula, 30-day aggregate & Long-established & -- & \ref{sec:m21} \\
Supervisory run-off, adverse path & Recent published & -- & \ref{sec:m23} \\
Supervisory run-off, extreme path & Recent published & -- & \ref{sec:m23} \\
Reverse-stress run-off multiplier with a refusal rule & Proposed in this work & -- & \ref{sec:bs11} \\
Daily compounding calibrated on the rows before 2023 & Reference & -- & -- \\
\bottomrule
\end{longtable}
}

The document adds a note on one of these methods. Reverse-stress run-off multiplier with a refusal rule: the multiplier is a threshold on a base, not a forecast: the score shown is the one of its calibrated fast-run base, and the comparators carry no score on the multiplier.

{\footnotesize\setlength{\tabcolsep}{3pt}
% bm-table Q10 normal cols=share/run;ratio/run;light/run;survival/run
\begin{longtable}{>{\raggedright\arraybackslash}p{0.25\textwidth}>{\raggedleft\arraybackslash}p{0.14\textwidth}>{\raggedleft\arraybackslash}p{0.14\textwidth}>{\raggedleft\arraybackslash}p{0.14\textwidth}>{\raggedleft\arraybackslash}p{0.14\textwidth}}
\caption{Reverse liquidity stress: Share of deposits lost in the two-day event horizon (share; nearer 0.85 is better); Ratio to the observed share (ratio; nearer 1 is better); Light against the observed run (light; green is best, then amber, then red); Survival period of the run-off ladder (days; context, not ranked), for each method, on the run The labelled run. The best value of a column by the declared rule is in bold, ties all in bold; a dash marks a cell without a value, and the notes below the table state why.}\label{tab:bm-q10-1}\\
\toprule Method & Share of deposits lost in the two-day event horizon & Ratio to the observed share & Light against the observed run & Survival period of the run-off ladder \\ \midrule \endfirsthead
\toprule Method & Share of deposits lost in the two-day event horizon & Ratio to the observed share & Light against the observed run & Survival period of the run-off ladder \\ \midrule \endhead
\multicolumn{5}{l}{\emph{Long-established}} \\
Liquidity coverage ratio formula, 30-day aggregate & -- & \textbf{0.506} & \textbf{amber} & -- \\ % bm-row lcr
\multicolumn{5}{l}{\emph{Recent published}} \\
Supervisory run-off, adverse path & 0.091 & 0.108 & red & 62 \\ % bm-row list-adverse
Supervisory run-off, extreme path & 0.098 & 0.115 & red & 41 \\ % bm-row list-extreme
\multicolumn{5}{l}{\emph{Proposed in this work}} \\
Reverse-stress run-off multiplier with a refusal rule & 0.171 & 0.201 & red & -- \\ % bm-row multiplier
\multicolumn{5}{l}{\emph{Reference}} \\
Daily compounding calibrated on the rows before 2023 & \textbf{0.171} & 0.201 & red & -- \\ % bm-row fast-run
\bottomrule
\end{longtable}
}

A dash in Table~\ref{tab:bm-q10-1} is a cell the product does not compute, for the stated reason: the formula has no daily path, so it has no outflow at the two-day horizon (Liquidity coverage ratio formula, 30-day aggregate); the base survives its horizon (Daily compounding calibrated on the rows before 2023). For Liquidity coverage ratio formula, 30-day aggregate and Reverse-stress run-off multiplier with a refusal rule the document holds no entry in some columns of Table~\ref{tab:bm-q10-1}; that is an entry the document does not hold, not a value of zero, and nothing is claimed for it.

{\footnotesize\setlength{\tabcolsep}{3pt}
% bm-table Q10 normal cols=multiplier/b-list-adverse;multiplier/b-list-extreme;multiplier/b-fast-run
\begin{longtable}{>{\raggedright\arraybackslash}p{0.25\textwidth}>{\raggedleft\arraybackslash}p{0.14\textwidth}>{\raggedleft\arraybackslash}p{0.14\textwidth}>{\raggedleft\arraybackslash}p{0.14\textwidth}}
\caption{Reverse liquidity stress: Run-off multiplier at the 30-day target (multiplier; context, not ranked), for each method, on the run The labelled run. The best value of a column by the declared rule is in bold, ties all in bold; a dash marks a cell without a value, and the notes below the table state why.}\label{tab:bm-q10-2}\\
\toprule Method & Base: adverse run-off & Base: extreme run-off & Base: calibrated fast run \\ \midrule \endfirsthead
\toprule Method & Base: adverse run-off & Base: extreme run-off & Base: calibrated fast run \\ \midrule \endhead
\multicolumn{4}{l}{\emph{Proposed in this work}} \\
Reverse-stress run-off multiplier with a refusal rule & 2.118 & 1.383 & -- \\ % bm-row multiplier
\bottomrule
\end{longtable}
}

A dash in Table~\ref{tab:bm-q10-2} is a cell the product does not compute, for the stated reason: the base is refused: it falls below the 30-day target before any multiplier applies (Reverse-stress run-off multiplier with a refusal rule).

The document reads the tables of this question by the declared rule of each metric. Share of deposits lost in the two-day event horizon (nearer 0.85 is better): best in Daily compounding calibrated on the rows before 2023 1 of 1. Ratio to the observed share (nearer 1 is better): best in Liquidity coverage ratio formula, 30-day aggregate 1 of 1. Light against the observed run (green is best, then amber, red): best in Liquidity coverage ratio formula, 30-day aggregate 1 of 1. Reverse-stress run-off multiplier with a refusal rule, proposed in this work, is not ranked: the multiplier is a threshold on a base, not a forecast: the score shown is the one of its calibrated fast-run base, and the comparators carry no score on the multiplier. A null cell is never ranked; each states why it is null.
